\documentclass[journal]{IEEEtran}
\usepackage{cite}
\usepackage{amsmath,amssymb,amsfonts}
\usepackage{algorithmic}
\usepackage{graphicx}
\usepackage{textcomp}
\usepackage{xcolor}
\usepackage{url}
\usepackage[colorlinks=true, urlcolor=magenta]{hyperref}

\begin{document}

\title{Week 6 Assignment: Image Denoising with FFDNet and Comparison with DnCNN/BM3D/ISTA/FISTA/ADMM}

\author{Author Name}

\maketitle

\begin{abstract}
This report provides a comparative analysis of six distinct image denoising algorithms: Block-Matching and 3D filtering (BM3D), three iterative optimization methods (ISTA, FISTA, ADMM) solving the Total Variation (TV) model, and two deep learning approaches (DnCNN and FFDNet). Experiments demonstrate that deep learning methods maintain exceptional recovery quality while significantly accelerating inference speed. FFDNet, in particular, exhibits outstanding inference efficiency and superior structural preservation, significantly outperforming traditional optimization frameworks. The source code and reproducible experiments are available at: \url{https://github.com/WRw5w/exp6}.
\end{abstract}

\section{Introduction}
Image denoising aims to recover a clean image $x$ from a noisy observation $y = x + v$, where $v$ is typically assumed to be Additive White Gaussian Noise (AWGN). Optimization methods iteratively solve an objective function based on hand-crafted physical priors, which is mathematically rigorous but computationally expensive. In contrast, deep learning methods utilize large datasets to learn a direct mapping from noisy to clean images, offering rapid inference. While DnCNN achieves impressive performance at a single noise level through simple convolutional stacking and residual learning, FFDNet introduces paradigm-shifting efficiency. FFDNet achieves this by operating on downsampled sub-images rather than processing the original image directly, which significantly accelerates inference speed and expands the receptive field. Furthermore, unlike DnCNN which requires a separate model for each noise level, FFDNet takes a noise level map as input, allowing a single network to flexibly handle spatially variant and unknown noises without retraining \cite{zhang2018ffdnet}.

\section{Methods}

\subsection{Block-Matching and 3D Filtering (BM3D)}
BM3D is a state-of-the-art traditional algorithm relying on non-local patch similarity. The core concept involves grouping similar 2D image fragments into 3D data arrays, applying a 3D collaborative filter (e.g., using DCT or wavelet transforms), and then aggregating the estimates back to their original positions.

\subsection{Total Variation (TV) Minimization}
Optimization-based methods minimize the following Total Variation objective function:
\begin{equation}
\min_x \frac{1}{2}\|x-y\|_2^2 + \lambda \|\nabla x\|_1
\end{equation}
Typical solvers include \textbf{ISTA} (Iterative Shrinkage-Thresholding Algorithm), which applies a forward gradient descent step followed by a proximal mapping on the TV penalty; \textbf{FISTA}, a fast variant of ISTA that introduces a Nesterov-style momentum term accelerating the convergence rate from $\mathcal{O}(1/k)$ to $\mathcal{O}(1/k^2)$; and \textbf{ADMM} (Alternating Direction Method of Multipliers), which tackles the problem by splitting variables and solving subproblems via an augmented Lagrangian, yielding faster empirical convergence.

\begin{figure*}[t]
\centering
\includegraphics[width=0.9\textwidth]{outputs/denoising_visual_comparison.jpg}
\caption{Visual comparison of denoising methods on the test image. FFDNet and DnCNN achieve the highest structural similarity to the ground truth.}
\label{fig:visual}
\end{figure*}



\subsection{Denoising Convolutional Neural Network (DnCNN)}
DnCNN performs denoising through a deep fully convolutional architecture utilizing residual learning. Instead of directly outputting the denoised image $x$, the network $\mathcal{R}(\cdot)$ learns to predict the residual noise $v$:
\begin{equation}
\mathcal{R}(y) \approx v \implies \hat{x} = y - \mathcal{R}(y)
\end{equation}

\subsection{Fast and Flexible Denoising Network (FFDNet)}
FFDNet inherits the power of CNNs but completely shifts the data pathway. It employs a reversible downsampling operator (e.g. PixelShuffle-like strategy) to rearrange the input image and an auxiliary noise variance map into four orthogonal sub-image channels. It applies non-linear convolutions in this lower-resolution domain before upsampling it back to the original dimensions. This decoupling mechanism drastically reduces runtime complexity while providing robust preservation of high-frequency structures.

Furthermore, unlike traditional networks that require retraining for different noise levels, FFDNet accepts a noise level map as a tunable input map. This flexibility enables it to handle spatially variant noise and strike an optimal balance between visual quality, generalizability, and computational expediency during deployment.

\begin{figure}[htbp]
\centering
\includegraphics[width=\columnwidth]{ffdnet_arch.pdf}
\caption{Schematic diagram of the FFDNet architecture. It consists of a downsampling layer, non-linear convolutional layers, and an upsampling layer.}
\label{fig:ffdnet_arch}
\end{figure}

\section{Experiments}

\textbf{Setup:} We evaluated the algorithms on a standard \texttt{camera} image contaminated with AWGN ($\sigma = 25/255$). For the optimization methods (ISTA, FISTA, ADMM), we tuned the regularization parameter $\lambda=0.1$ and executed 100 iterations. For BM3D, we used the analytical baseline. For DnCNN and FFDNet, we utilized the pre-trained \texttt{dncnn\_25.pth} and \texttt{ffdnet\_gray.pth} models.

\textbf{Results:}
The quantitative performance measured in Peak Signal-to-Noise Ratio (PSNR) and Structural Similarity Index (SSIM) is summarized in Table \ref{tab:metrics}. 

\begin{table}[htbp]
\caption{Quantitative Comparison of Denoising Methods ($\sigma=25$)}
\label{tab:metrics}
\centering
\begin{tabular}{|l|c|c|c|}
\hline
\textbf{Method} & \textbf{PSNR (dB)} & \textbf{SSIM} & \textbf{Time (s)} \\
\hline
Noisy & 20.61 & 0.3012 & N/A \\
ADMM & 27.75 & 0.6438 & 0.3017 \\
ISTA & 28.52 & 0.7216 & 2.3575 \\
FISTA & 28.53 & 0.7221 & 2.1701 \\
BM3D & 29.70 & 0.7963 & 12.0372 \\
DnCNN & \textbf{29.99} & 0.8148 & 0.7471 \\
FFDNet & 29.98 & \textbf{0.8193} & \textbf{0.0555} \\
\hline
\end{tabular}
\end{table}

As shown in Figure \ref{fig:visual}, deep learning methods effectively remove noise while preserving sharp edges over traditional approaches. Furthermore, the execution speed and absolute performance metrics are illustrated in Figure \ref{fig:charts}. FFDNet demonstrates an unparalleled latency of merely $0.0555$s, completely altering the applicability limits of complex network inferences.



\begin{figure}[htbp]
\centering
\includegraphics[width=\columnwidth]{outputs/performance_charts.jpg}
\caption{PSNR/SSIM evaluation and logarithmic execution time scaling across the tested algorithms.}
\label{fig:charts}
\end{figure}

\section{Conclusion}
This study evaluated classical and learning-based denoising frameworks, taking a thorough examination of FFDNet's structural and empirical capabilities. The results confirmed that while optimization techniques like BM3D provide excellent recovery benchmarks, advanced dimensionality-driven network architectures, such as FFDNet, utilize coordinate rearrangements to achieve supreme compression of operational cost globally. It sets an industry-leading standard for achieving superior image fidelity, exceptional generalization, and practically real-time applicability in modern scenarios.

\begin{thebibliography}{00}
\bibitem{zhang2018ffdnet} K. Zhang, W. Zuo and L. Zhang, ``FFDNet: Toward a Fast and Flexible Solution for CNN-Based Image Denoising,'' in \textit{IEEE Transactions on Image Processing}, vol. 27, no. 9, pp. 4608-4622, Sept. 2018.
\end{thebibliography}

\end{document}
